\section{Introduction}
\label{sec:intro}

Contrastive vision--language models are often described as behaving like bags of words: they
retrieve images that contain the right objects and attributes but do not reliably tell apart
``a red cube and a blue sphere'' from ``a blue cube and a red sphere''
\citep{yuksekgonul2023bow,thrush2022winoground,lewis2024clipbind}. A common response is to train with
minimal-change counterfactuals, either as hard negatives \citep{yuksekgonul2023bow,zhang2024contrasting,
zhang2024countercurate,awal2024vismin} or with additional structure such as similarity-level
equivariance \citep{wang2023eqsim} or object-centric binding modules \citep{assouel2025occlip}.
Whether such training yields binding that transfers to attribute--object combinations or edit
compositions absent from training is an empirical question, and answering it requires an
evaluation that cannot be passed without binding.

This paper is about that evaluation rather than about a new method. Building it carefully exposed
three failure modes that are easy to miss. First, the input can carry the answer: in our own first
version, the ``image'' seen by the model was one token per object assembled from the scene graph,
so each token already bound its object's shape, colour and material. A head trained on such
tokens can only show how well the text side is aligned, not whether binding is learned from
visual input. Second, a metric can be solvable without binding: when minimal-change groups include
edits that change which colours or shapes are present, a binding-blind content channel scores well
above chance \citep[cf.][]{hsieh2023sugarcrepe}. Third, blind controls can look reassuring for
structural reasons: on $2{\times}2$ groups a single-modality scorer ties on every comparison, so
its group score is zero by construction, and on a pair metric that pairs every caption with one
matching and one non-matching image its AUC is exactly $0.5$. Neither number says anything about
whether the benchmark contains exploitable artefacts \citep[cf.][]{lin2024languagepriors}.

We therefore use a scene-graph generator in which every group of two images and two captions is
validated by an oracle, render the scenes to RGB, and give the model only pixels and caption
strings. Splits are made at the level of edit orbits, so no content multiset seen in training
reappears in development or test data, and the endpoint is the subset of \emph{binding-necessary}
groups, whose two scenes have identical object and attribute multisets. Within this setting we
examine the simplest reasonable learner: a small head on frozen CLIP ViT-B/32 tokens
\citep{radford2021clip,ilharco2021openclip}, trained with a hard-negative contrastive loss. We also
state what a cross-group edit-consistency objective adds to that loss for the inner-product scorers
used here.

Our current observations are narrow and mostly negative. On a small development panel the head
fits its training groups but performs at or below chance on binding-necessary development groups,
while a binding-blind content channel reproduces most of its score on the full panel. Oracle
object tokens used as a control give higher, but still low, development scores, so we cannot yet
separate limits of the frozen encoder from limits of the training data or the head. On a larger
panel we also ran a paired comparison of the hard-negative loss with and without the edit-consistency
objective under study. Neither arm fitted its training groups under the fixed update budget, and
training was unstable, so that comparison is unresolved rather than negative or positive. The
held-out test splits have not been opened.

\paragraph{Contributions.} (i) A description of an oracle-validated, minimal-change evaluation
setting for compositional binding with pixel-only model inputs, orbit-level splits and a
metadata-defined binding-necessary endpoint (Section~\ref{sec:setting}). None of these ingredients is new
on its own; we document how they interact. (ii) Derivations of chance levels and tie handling for
each group metric, and the structural properties of single-modality controls (Section~\ref{sec:setting},
Appendices~\ref{app:chance} and~\ref{app:blind}). (iii) The standard bilinear identity that relates hard-negative
training, the GroupMatch metric \citep{zhu2025ttm} and within-group edit alignment, which isolates
the cross-group part of the edit-consistency objective under study (Section~\ref{sec:method}).
(iv) Development-panel observations for a hard-negative baseline on frozen CLIP tokens, with
content-only, shuffle, blind and oracle-token controls. We add a pre-registered paired comparison
with the edit-consistency objective, whose outcome under the pre-fixed rules is ``optimization
unresolved'' (Section~\ref{sec:experiments}). (v) A metadata analysis of what orbit-level
splitting does to the generated data (Appendix~\ref{app:splits}).
